\scanpage{40}
\deleted{Il reste à définir $\unsure{b|B}$.}
Il faut expliciter $K(2)$, \deleted{\unsure{comme}} en choisissant une origine
\[
(\unsure{50|p0})\qquad\theta_\infty=
(\underbrace{\sqrt[n]{2}}_{\theta_n})_{n\in\mathbb N^*}
\]
donc $K(2)$ est identifié à $T$ donc à $\widehat{\mathbb Z}$, et
\deleted{$\unsure{A}$}
\[
(51)\qquad A\simeq\operatorname{Aff}(1,\widehat{\mathbb Z})
\simeq\widehat{\mathbb Z}^{*}\mathbin{\underset{\scriptscriptstyle1/2}{.}}\widehat{\mathbb Z}
\]
opérant sur $K(2)\simeq\widehat{\mathbb Z}$ de façon tautologique
\[
(52)\qquad
(\underset{\substack{\rotatebox{90}{$\in$}\\\widehat{\mathbb Z}^{*}}}{a},
\underset{\substack{\rotatebox{90}{$\in$}\\\widehat{\mathbb Z}}}{b})
(\underset{K(2)}x)=ax+b
\]
On a donc
\[
(53)\qquad K(\tfrac12)\simeq\widehat{\mathbb Z}
\]
mais opérations
\[
(54)\qquad
(\underset{\substack{\rotatebox{90}{$\in$}\\\widehat{\mathbb Z}^{*}}}{a},
\underset{\substack{\rotatebox{90}{$\in$}\\\widehat{\mathbb Z}}}{b})
(\underset{K(\frac12)}x)=ax-b
\]
Ceci permet d'expliciter entièrement $E(\unsure{\beta_*|\beta_\alpha})$
en plus de $E(\omega)\ldots$

Il reste \inserted{enfin} à expliciter les hom de transport parallèle
\diagram{\[
\begin{tikzcd}[ampersand replacement=\&,row sep=small,column sep=large]
\&E_0\&\&\&E'_0=E_2\\
E(\alpha)\arrow[ur]\arrow[r]\arrow[dr]\&E_1\&\text{et}\&
E(\alpha)\arrow[ur,"\sim"]\arrow[r,"\sim"]\arrow[dr,"\sim"']\&E'_1=E_{-1}\\
\&E_\infty\&\&\&E'_\infty\simeq E_{1/2}
\end{tikzcd}
\]}
\inserted{et
\diagram{\[
\begin{tikzcd}[ampersand replacement=\&,row sep=small,column sep=small]
\&E_1\&\&E_\infty\&\&E_0\\
E'_0\arrow[ur]\arrow[r]\&E_\infty\&
E'_1\arrow[ur]\arrow[r]\&E_0\&
E'_\infty\arrow[ur]\arrow[r]\&E_1
\end{tikzcd}
\]}}
d'où se déduit les hom. correspondants de $E(\overline\alpha)$
en conjuguant par la conjugaison complexe. Comme ce sont des hom
\margin{?}
% archive foliation: 39
